\section{Experiments}
\label{sec:experiments}
% Generated from audited RTX 3090 evidence: e30a75542a7f25f48d3218c8519fab993903d4f11d102973146a4fa41e8e9274
\newcommand{\QKSieveMhaAttentionRows}{%
8K & 0.1691 & 0.1327 & 0.1560 & 1.27$\times$ & 1.08$\times$ \\%
16K & 0.3173 & 0.1900 & 0.2250 & 1.67$\times$ & 1.41$\times$ \\%
32K & 0.6163 & 0.2324 & 0.2945 & 2.65$\times$ & 2.09$\times$ \\%
64K & 1.2027 & 0.2646 & 0.3576 & 4.55$\times$ & 3.36$\times$ \\%
128K & 2.3708 & 0.3722 & 0.5757 & 6.37$\times$ & 4.12$\times$ \\%
}

\newcommand{\QKSieveMhaFierRows}{%
8K & 0.1816 & 0.1327 & 0.1560 & 1.37$\times$ & 1.16$\times$ \\%
16K & 0.2951 & 0.1900 & 0.2250 & 1.55$\times$ & 1.31$\times$ \\%
32K & 0.4647 & 0.2324 & 0.2945 & 2.00$\times$ & 1.58$\times$ \\%
64K & 0.7329 & 0.2646 & 0.3576 & 2.77$\times$ & 2.05$\times$ \\%
128K & 1.3322 & 0.3722 & 0.5757 & 3.58$\times$ & 2.31$\times$ \\%
}

\newcommand{\QKSieveMhaDecodeRows}{%
32K & 84.18 & 51.56 & 63.97 & 1.63$\times$ & 1.32$\times$ \\%
64K & 144.75 & 52.97 & 65.25 & 2.73$\times$ & 2.22$\times$ \\%
128K & 268.17 & 55.41 & 67.38 & 4.84$\times$ & 3.98$\times$ \\%
}

\newcommand{\QKSieveMhaBuildBreakEvenRows}{%
32K & 0.768 & 1.375 & 24 & 69 & 1.01$\times$ & 0.81$\times$ \\%
64K & 0.774 & 1.488 & 9 & 19 & 1.64$\times$ & 1.31$\times$ \\%
128K & 0.743 & 1.839 & 4 & 10 & 2.75$\times$ & 2.15$\times$ \\%
}

\newcommand{\QKSievePersistentRows}{%
32K & 0.702$\times$ & 1.322$\times$ & 1.082$\times$ & 1.319$\times$ & 1.759 \\%
64K & 1.125$\times$ & 2.221$\times$ & 1.785$\times$ & 2.211$\times$ & 1.998 \\%
}


We test equal-byte retrieval geometry and allocation, frozen downstream
quality across tasks and models, and whole-model crossover after all
overhead.

\subsection{Protocol}

Both main profiles use $\lambda=0.75$, stride-32 Key samples, eight
prompt-tail Queries, 240-bit plain-scale indices, complete packed-index
qMSE allocation, and \cref{eq:budget-intro}.  The reference materializes all
proxy scores and applies top-$k$; the optimized path replaces it with the
capped sampled-quantile scan in \cref{eq:frozen-quantile-samples}.  Its
request-local transform, allocation, exact-KV consumer, and budget are
unchanged.  We register Llama-3.1-8B-Instruct,
Qwen3-4B-Instruct, and
Mistral-7B-Instruct-v0.3, complete English LongBench, RULER from 4K to
128K, and paired held-out-position PPL.

Synchronized batch-one timing reports complete attention, whole-model steady
forward, a directly timed decode horizon, and a directly timed request that
includes prefill and initialization.  Every auxiliary stage is also invoked
and timed independently for diagnosis, but no complete-path latency is
constructed by summing those measurements.  RTX~3090 numbers are audited
diagnostics.
\IfFileExists{data/generated/qksieve_h100_tables.tex}{%
Matched H100 measurements are reported in
\cref{tab:h100-attention,tab:h100-decode,tab:h100-requests}.}{%
Matched H100 measurements remain required.}

Our claimed scope is decode-time attention over an already materialized,
reusable KV cache, rather than sparse prefill.  We therefore keep cold
single-use, same-basis shared-prefix branching, and append-only timing as
separate contracts.  An audited 32/64K slice is reported below.
\IfFileExists{data/generated/qksieve_h100_tables.tex}{%
The matched 64/128K H100 lifecycle grid is reported in
\cref{tab:h100-requests}.}{%
The remaining H100/128K grid is preregistered in
\cref{app:persistent-kv-protocol}.}
Cold request
latency remains a limitation rather than being mixed with the warm-path claim.

\IfFileExists{data/generated/qksieve_quality_tables.tex}{%
% Generated from frozen evidence: longbench=382f0f892e8d1e098fc94b5d966a6b759d018ff267952e0cfa71127ee5dd7a1f; ruler=bf945639ecd056b0dc0e7e6411be5c081a0fcaea7821caffa41028eae77c82b1; multimodel=9ceadaace51808989a222df6669ca5261e233c8ebd21f67b3f2bcd9851b8bf30
\begin{table}[t]
\caption{Task quality for frozen QKSieve-Robust. Complete LongBench, RULER, and each cross-model screen contain 3,750, 650, and 160 strict pairs; intervals bootstrap tasks or task--length cells.}
\label{tab:quality-main}
\centering
\small
\resizebox{\columnwidth}{!}{%
\begin{tabular}{@{}lrrrr@{}}
\toprule
Evaluation / model & Full & QKSieve & Retention & Paired 95\% CI \\
\midrule
Full LB / Llama-3.1-8B & 0.4590 & 0.4587 & 99.93\% & [99.54\%, 100.21\%] \\
RULER / Llama-3.1-8B & 0.8416 & 0.8444 & 100.33\% & [99.18\%, 101.75\%] \\
\midrule
LB screen / Llama-3.1-8B & 0.4264 & 0.4208 & 98.68\% & [96.39\%, 100.50\%] \\
LB screen / Qwen3-4B & 0.3972 & 0.3980 & 100.21\% & [98.91\%, 101.72\%] \\
LB screen / Mistral-7B & 0.4216 & 0.4153 & 98.49\% & [95.08\%, 100.89\%] \\
\bottomrule
\end{tabular}%
}
\end{table}

}{%
The complete LongBench table uses the request-local reference profile.  A
strict smaller screen then runs the sampled-quantile path used by the CUDA
timing experiment.

\begin{table}[t]
\caption{Complete reference-profile LongBench result.  The optimized
deployment profile is evaluated separately.}
\label{tab:longbench-main}
\centering
\begin{tabular}{lrrrr}
\toprule
Model & Full & \method & Retention & Paired 95\% CI \\
\midrule
Llama-3.1-8B & .459398 & .458852 & 99.881\% &
[99.424\%, 100.347\%] \\
\bottomrule
\end{tabular}
\end{table}

This covers 3,750 strict pairs and all 16 tasks; per-task scores, including
the 97.615\% minimum retention, are in \cref{app:longbench}.

On 160 strict same-sample pairs, an earlier length-gated Fast/Robust path
scores .453676 versus .453019 for Full, retaining 100.145\%.  A paired-sample
bootstrap gives [97.915\%, 102.051\%].  This closes the
reference/deployment code-path mismatch at screening scale, but the interval
and localized Qasper/MultiNews ratios of 84.28/92.85\% preclude an
``equivalent quality'' or superiority claim.  We retain the complete
reference result as the primary task-quality evidence.

A separate targeted RULER audit covers 13 tasks at 8/16/32/64K with
2/2/2/1 examples per task and length.  Across 91 strict pairs, Full and
\method score .884776 and .884135, retaining 99.928\% with a paired-bootstrap
interval of [99.559\%, 100.294\%].  This small audit checks task-family and
length failure modes; it is not presented as a formal full-size RULER table,
and its quality-harness timing is excluded from systems claims.
}

\IfFileExists{data/generated/qksieve_quality_tables.tex}{%
PassageCount is the only LongBench task below 95\% relative retention:
.0991 for Full versus .0916 for Robust (92.43\%).  We report this case
explicitly because its small Full denominator makes the relative ratio look
larger than the absolute $.0075$ decrease.}{ }

\subsection{Index quality at equal rate}

At 240 index bits and 1,280 exact tokens over 12 held-out 32K windows,
QK balancing raises Full-token top-1 agreement by 1.95 points and reduces
KL by 60\% for .81\% latency (\cref{tab:key-only-causal}).  Key-MSE matches
logit-MSE downstream and is 2.46\% faster
(\cref{tab:allocation-causal,tab:allocation-deploy}).

At 1\% active KV, automatic QK-MSE reaches 70.48\% exact-set and 96.36\%
oracle-mass recall; its matched-rate result is close to fixed 4-4-2-1.

At 1/2/4\% active KV, paired attention-mass gains over FIER are
7.08/5.68/4.35 points with intervals excluding zero \citep{wang2025fier}.
Across 59,200 conditions, matched-rate RaBitQ recall/mass rises from
53.88/71.32\% to 71.82/76.79\% (\cref{app:mechanism-results}).

On a strict 160-sample LongBench screening set, the Full macro score is
.45302.  At the same active-token schedule, the gated ablation, the
complete-top-$k$ \method reference, an audited FIER RTN-1 g32 reference,
Quest-P16, a RaBitQ formula reference, SparQ-R32 selector-only, the complete
SparQ formula, and a published-equation UNIQUE-P8 reference retain
100.15/99.79/100.01/92.22/100.01/99.62/99.95/93.19\%, respectively.  The
absolute macro difference between the gated ablation and FIER is .00061; this small
screen does not support a superiority claim.  These rows compare selector
formulas and downstream quality, not official kernel speed.  On the same
160 pairs, BinaryPC's released Llama projection retains 99.25\%.  The gated ablation
invokes neither a task router nor Full fallback, but ten examples per task
remain too few for a superiority claim.

\subsection{Length scaling}

\begin{table}[t]
\caption{Frozen native-MHA attention paths on one RTX~3090.  Full is native
FP16 SDPA without KV replication.  Fast includes the selector and exact sparse
attention; Robust additionally includes the rank-16 INT4 Value-tail scan and
consumer.  Index construction is excluded from all rows.}
\label{tab:mha-main}
\centering
\small
\begin{tabular}{rrrrrr}
\toprule
History & Full & Fast & Robust & Fast sp. & Robust sp. \\
\midrule
\QKSieveMhaAttentionRows
\bottomrule
\end{tabular}
\end{table}

The frozen timing shape is batch one, 32 Query heads, 32 KV heads, and
$d=128$, with resident FP16 K/V.  CUDA-event measurements use 10 warmup and
40 timed iterations; each row is the median of three independent seeds.
The sparse path includes fused Query
projection/quantization, mixed-bit scanning, sampled-threshold candidate
compaction, and exact QK--softmax--AV.  Fast and Robust have exactly equal
thresholds, candidate counts, and candidate sets in all 15 seed--length
comparisons.  Their only difference is omitted-Value estimation.  The
allocated Key-only and Key-plus-Value auxiliary tensors occupy approximately
5.47\% and 7.09\% of FP16 K+V in this kernel; the method-level hard rate caps,
including amortized metadata, are 5.86\% and 7.47\%.

\begin{table}[t]
\caption{Controlled comparison with the audited packed FIER RTN-1 g32
implementation.  FIER and Fast use the same active-token schedule and exact
sparse-attention consumer; Robust additionally estimates the omitted Value
tail.  Ratios above one favor QKSieve.}
\label{tab:mha-fier-samepath}
\centering
\small
\begin{tabular}{rrrrrr}
\toprule
History & FIER & Fast & Robust & Fast/FIER & Robust/FIER \\
\midrule
\QKSieveMhaFierRows
\bottomrule
\end{tabular}
\end{table}

\Cref{tab:mha-fier-samepath} changes neither model tensors nor the final exact
K/V consumer between FIER and Fast, so it isolates packed-selector cost.  The
Robust column is a stricter systems comparison because it also pays for the
Value-tail estimator.  This is a controlled implementation comparison, not a
claim about an unmeasured vendor-optimized FIER stack.

\Cref{tab:mha-main} is an attention-subsystem result, not an end-to-end decode
claim.  It excludes prefill and historical index construction.  We therefore
measure the same Fast and Robust paths inside a real MHA model separately.

\begin{table}[t]
\caption{Steady greedy decode on Yarn-Llama-2-7B-128K with resident exact K/V.
Each run generates 64 tokens and reports the final 48.  Values are directly
timed whole-model milliseconds per token.}
\label{tab:mha-decode}
\centering
\small
\begin{tabular}{rrrrrr}
\toprule
History & Full & Fast & Robust & Fast sp. & Robust sp. \\
\midrule
\QKSieveMhaDecodeRows
\bottomrule
\end{tabular}
\end{table}

Fast/Robust index construction takes .768/1.375, .774/1.488, and
.743/1.839 seconds at 32/64/128K, giving break-even lengths of 24/69, 9/19,
and 4/10 tokens and directly timed 64-token online speedups of 1.01/.81,
1.64/1.31, and 2.75/2.15$\times$.  Dense prefill and deployment-time CUDA JIT
are excluded; see \cref{tab:mha-build-break-even}.

\begin{table}[t]
\caption{Directly timed persistent-cache lifecycle on the same MHA model and
RTX~3090 host.  Cold-index includes the Robust build and one 32-token branch
but excludes dense prefill; cold-E2E includes prefill.  Warm reuses the index;
four-branch amortization shares one build over 128 generated tokens.  Entries
are medians and 95\% process-repetition intervals from three independent
processes on one fixed deterministic workload.}
\label{tab:persistent-kv-main}
\centering
\resizebox{\linewidth}{!}{
\begin{tabular}{rrrrrrr}
\toprule
History & Cold-index & Cold-E2E & Warm & 4-branch & Append-only & Build (s) \\
\midrule
\QKSievePersistentRows
\bottomrule
\end{tabular}}
\end{table}

For 32/64K, median Full/QKSieve warm latency is
\QKSievePersistentWarmLatencyText\ ms/token; median Robust construction is
\QKSievePersistentBuildText\ seconds.  The matched
$\lceil I/(D_f-D_s)\rceil$ break-even is
\QKSievePersistentBreakEvenText\ generated tokens.  The 32K cold-E2E result
is therefore honestly below $1\times$, whereas 64K is approximately break-even
for a single 32-token request and benefits substantially once the prefix is
reused.

The lifecycle audit covers all 32 layers, six snapshots, and five rewinds,
with stable index pointers and deterministic replay.  Reuse is valid while
request-local Query statistics remain fixed; a new prompt requires rebuilding
the request-local index.  Here append-only means one uninterrupted 128-token
autoregressive continuation whose exact K/V and auxiliary index entries grow
one token at a time.  It does not include delta prefill for a new user turn.

\IfFileExists{data/generated/qksieve_h100_tables.tex}{%
\input{data/generated/qksieve_h100_tables.tex}}{}

At the native 128K boundary, a corrected same-seed six-topic audit reverses
our initial conclusion.  QKSieve-Fast retains only 96.319\% pooled PPL
quality (96.373\% case macro) and falls to 91.100\% on politics.  Exact-QK
top-1,280 also retains only 94.74\% on a medicine attribution window, and
doubling its budget reaches 96.24\%; the failure is therefore not solely
low-bit Key misranking.  Adding 128 or 256 recent tokens changes the weak-case
macro by less than .1 point.

QKSieve-Robust rank-16 INT4 with $\alpha=.5$ is frozen on the six same-seed
cases and evaluated on six new seeds.  Across all 12 cases it retains
99.805\% pooled PPL quality with a case-bootstrap 95\% interval of
[99.100\%, 100.499\%], a 99.812\% case macro, and a 96.980\% worst case.
The total per-token auxiliary rate is 7.471\%.  Its earlier GQA-4 subsystem
and whole-model timings are retained in \cref{app:system-diagnostics}; they
are not combined with the native-MHA speed claim in \cref{tab:mha-main}.

Length-gated ablations and the 256K coordinate-extrapolation failure are
reported in \cref{app:system-diagnostics}; neither changes the frozen
always-on Robust contract or its registered 4K--128K scope.
